\documentclass[exercice]{thedocument}%
\startdatakeys
\source{}
\BO{2026}
\cycle{4}
\competencesDuSocle{CA}
\objectifsApprentissage{C4A1bm,C4A1bb,C4A1eb}
\automatismes{C4A1ch,C4C2aa}
\enddatakeys
\begin{document}%
En utilisant les quatre nombres (-7) ; (+15) ; (-7,6) et (+3,9) une fois
chacun et des signes d'addition et/ou de soustraction, écrire une
expression numérique dont la valeur est~:
\begin{enumerate}
    \item le plus grand nombre possible~;\vskip50pt
    \item le plus petit nombre possible~;\vskip50pt
    \item le nombre le plus proche de zéro.\vskip50pt
\end{enumerate}
\end{document}
